% Bewegungsplanung im Konfigurationsraum: Pfad von q_I nach q_G im freien Raum.
% Eigene Darstellung in Anlehnung an LaValle2006 (Fig. 4.11, S. 71; Beleg LaValle2006.07).
% Stil: ba-grafik.tex. Einbinden: % Bewegungsplanung im Konfigurationsraum: Pfad von q_I nach q_G im freien Raum.
% Eigene Darstellung in Anlehnung an LaValle2006 (Fig. 4.11, S. 71; Beleg LaValle2006.07).
% Stil: ba-grafik.tex. Einbinden: % Bewegungsplanung im Konfigurationsraum: Pfad von q_I nach q_G im freien Raum.
% Eigene Darstellung in Anlehnung an LaValle2006 (Fig. 4.11, S. 71; Beleg LaValle2006.07).
% Stil: ba-grafik.tex. Einbinden: % Bewegungsplanung im Konfigurationsraum: Pfad von q_I nach q_G im freien Raum.
% Eigene Darstellung in Anlehnung an LaValle2006 (Fig. 4.11, S. 71; Beleg LaValle2006.07).
% Stil: ba-grafik.tex. Einbinden: \input{Bilder/tikz/konfigurationsraum} in einer figure-Umgebung.
\begin{tikzpicture}[ba]
  % Konfigurationsraum C
  \draw[ba normal, baGrau, fill=white] plot[smooth cycle, tension=0.8] coordinates {(0,0.4) (2.5,-0.2) (5.6,0.2) (7.2,1.6) (6.6,3.6) (3.8,4.2) (1.0,3.8) (-0.4,2.2)};
  % Hindernisbereiche C_obs
  \fill[baHellgrau] plot[smooth cycle, tension=0.9] coordinates {(2.2,1.2) (3.4,0.8) (3.9,1.9) (3.1,2.9) (2.3,2.4)};
  \fill[baHellgrau] plot[smooth cycle, tension=0.9] coordinates {(4.6,2.3) (5.7,1.9) (6.3,2.9) (5.4,3.6) (4.5,3.2)};
  \node at (3.05,1.85) {$C_\mathrm{obs}$};
  \node at (5.4,2.8) {$C_\mathrm{obs}$};
  % Pfad
  \coordinate (qI) at (0.7,1.6);
  \coordinate (qG) at (6.5,1.2);
  \draw[ba glied, line width=1.6pt] plot[smooth, tension=0.7] coordinates {(qI) (1.6,0.9) (3.2,0.35) (4.4,1.2) (5.4,1.3) (qG)};
  \fill[black] (qI) circle (1.2mm) node[above left] {$q_\mathrm{I}$};
  \fill[black] (qG) circle (1.2mm) node[above right] {$q_\mathrm{G}$};
  \node[anchor=west] at (1.0,3.3) {$C_\mathrm{free}$};
  \node[anchor=west, text=baGrau] at (7.3,3.4) {$C = C_\mathrm{free} \cup C_\mathrm{obs}$};
\end{tikzpicture}
 in einer figure-Umgebung.
\begin{tikzpicture}[ba]
  % Konfigurationsraum C
  \draw[ba normal, baGrau, fill=white] plot[smooth cycle, tension=0.8] coordinates {(0,0.4) (2.5,-0.2) (5.6,0.2) (7.2,1.6) (6.6,3.6) (3.8,4.2) (1.0,3.8) (-0.4,2.2)};
  % Hindernisbereiche C_obs
  \fill[baHellgrau] plot[smooth cycle, tension=0.9] coordinates {(2.2,1.2) (3.4,0.8) (3.9,1.9) (3.1,2.9) (2.3,2.4)};
  \fill[baHellgrau] plot[smooth cycle, tension=0.9] coordinates {(4.6,2.3) (5.7,1.9) (6.3,2.9) (5.4,3.6) (4.5,3.2)};
  \node at (3.05,1.85) {$C_\mathrm{obs}$};
  \node at (5.4,2.8) {$C_\mathrm{obs}$};
  % Pfad
  \coordinate (qI) at (0.7,1.6);
  \coordinate (qG) at (6.5,1.2);
  \draw[ba glied, line width=1.6pt] plot[smooth, tension=0.7] coordinates {(qI) (1.6,0.9) (3.2,0.35) (4.4,1.2) (5.4,1.3) (qG)};
  \fill[black] (qI) circle (1.2mm) node[above left] {$q_\mathrm{I}$};
  \fill[black] (qG) circle (1.2mm) node[above right] {$q_\mathrm{G}$};
  \node[anchor=west] at (1.0,3.3) {$C_\mathrm{free}$};
  \node[anchor=west, text=baGrau] at (7.3,3.4) {$C = C_\mathrm{free} \cup C_\mathrm{obs}$};
\end{tikzpicture}
 in einer figure-Umgebung.
\begin{tikzpicture}[ba]
  % Konfigurationsraum C
  \draw[ba normal, baGrau, fill=white] plot[smooth cycle, tension=0.8] coordinates {(0,0.4) (2.5,-0.2) (5.6,0.2) (7.2,1.6) (6.6,3.6) (3.8,4.2) (1.0,3.8) (-0.4,2.2)};
  % Hindernisbereiche C_obs
  \fill[baHellgrau] plot[smooth cycle, tension=0.9] coordinates {(2.2,1.2) (3.4,0.8) (3.9,1.9) (3.1,2.9) (2.3,2.4)};
  \fill[baHellgrau] plot[smooth cycle, tension=0.9] coordinates {(4.6,2.3) (5.7,1.9) (6.3,2.9) (5.4,3.6) (4.5,3.2)};
  \node at (3.05,1.85) {$C_\mathrm{obs}$};
  \node at (5.4,2.8) {$C_\mathrm{obs}$};
  % Pfad
  \coordinate (qI) at (0.7,1.6);
  \coordinate (qG) at (6.5,1.2);
  \draw[ba glied, line width=1.6pt] plot[smooth, tension=0.7] coordinates {(qI) (1.6,0.9) (3.2,0.35) (4.4,1.2) (5.4,1.3) (qG)};
  \fill[black] (qI) circle (1.2mm) node[above left] {$q_\mathrm{I}$};
  \fill[black] (qG) circle (1.2mm) node[above right] {$q_\mathrm{G}$};
  \node[anchor=west] at (1.0,3.3) {$C_\mathrm{free}$};
  \node[anchor=west, text=baGrau] at (7.3,3.4) {$C = C_\mathrm{free} \cup C_\mathrm{obs}$};
\end{tikzpicture}
 in einer figure-Umgebung.
\begin{tikzpicture}[ba]
  % Konfigurationsraum C
  \draw[ba normal, baGrau, fill=white] plot[smooth cycle, tension=0.8] coordinates {(0,0.4) (2.5,-0.2) (5.6,0.2) (7.2,1.6) (6.6,3.6) (3.8,4.2) (1.0,3.8) (-0.4,2.2)};
  % Hindernisbereiche C_obs
  \fill[baHellgrau] plot[smooth cycle, tension=0.9] coordinates {(2.2,1.2) (3.4,0.8) (3.9,1.9) (3.1,2.9) (2.3,2.4)};
  \fill[baHellgrau] plot[smooth cycle, tension=0.9] coordinates {(4.6,2.3) (5.7,1.9) (6.3,2.9) (5.4,3.6) (4.5,3.2)};
  \node at (3.05,1.85) {$C_\mathrm{obs}$};
  \node at (5.4,2.8) {$C_\mathrm{obs}$};
  % Pfad
  \coordinate (qI) at (0.7,1.6);
  \coordinate (qG) at (6.5,1.2);
  \draw[ba glied, line width=1.6pt] plot[smooth, tension=0.7] coordinates {(qI) (1.6,0.9) (3.2,0.35) (4.4,1.2) (5.4,1.3) (qG)};
  \fill[black] (qI) circle (1.2mm) node[above left] {$q_\mathrm{I}$};
  \fill[black] (qG) circle (1.2mm) node[above right] {$q_\mathrm{G}$};
  \node[anchor=west] at (1.0,3.3) {$C_\mathrm{free}$};
  \node[anchor=west, text=baGrau] at (7.3,3.4) {$C = C_\mathrm{free} \cup C_\mathrm{obs}$};
\end{tikzpicture}
